\documentclass[12pt, a4paper]{report}

% ── Packages ──────────────────────────────────────────────────────────────
\usepackage[utf8]{inputenc}
\usepackage[T1]{fontenc}
\usepackage{amsmath, amssymb, amsthm, mathtools, bm}
\usepackage{physics}
\usepackage{geometry}
\usepackage{hyperref}
\usepackage{cleveref}
\usepackage{enumitem}
\usepackage{titlesec}
\usepackage{fancyhdr}
\usepackage{booktabs}
\usepackage{array}
\usepackage{microtype}
\usepackage{xcolor}

\geometry{margin=1in}

% ── Hyperref setup ────────────────────────────────────────────────────────
\hypersetup{
    colorlinks=true,
    linkcolor=blue!70!black,
    citecolor=green!50!black,
    urlcolor=blue!60!black
}

% ── Theorem environments ──────────────────────────────────────────────────
\theoremstyle{definition}
\newtheorem{definition}{Definition}[chapter]
\newtheorem{example}{Example}[chapter]

\theoremstyle{plain}
\newtheorem{theorem}{Theorem}[chapter]
\newtheorem{proposition}{Proposition}[chapter]
\newtheorem{lemma}{Lemma}[chapter]

\theoremstyle{remark}
\newtheorem{remark}{Remark}[chapter]
\newtheorem{notation}{Notation}[chapter]
\newtheorem{warning}{Warning}[chapter]

% ── Custom commands ───────────────────────────────────────────────────────
\newcommand{\tildee}[1]{\tilde{e}_{#1}}
\newcommand{\tildeeps}[1]{\tilde{\epsilon}^{#1}}
\newcommand{\Rn}{\mathbb{R}^n}
\newcommand{\R}{\mathbb{R}}
\newcommand{\kd}{\delta}                    % Kronecker delta
\newcommand{\tp}{\otimes}                   % tensor product

% ── Header / Footer ──────────────────────────────────────────────────────
\pagestyle{fancy}
\fancyhf{}
\fancyhead[L]{\leftmark}
\fancyhead[R]{Tensors for Beginners}
\fancyfoot[C]{\thepage}

% ── Title ─────────────────────────────────────────────────────────────────
\title{%
    \vspace{-1cm}
    {\Huge\bfseries Tensors for Beginners}\\[0.8em]
    {\Large Lecture Notes}\\[0.4em]
    {\normalsize Based on the \emph{Tensors for Beginners} video series by eigenchris}
}
\author{}
\date{\today}

% ══════════════════════════════════════════════════════════════════════════
\begin{document}
\maketitle
\tableofcontents
\newpage

% ======================================================================
%  CHAPTER -1 : MOTIVATION
% ======================================================================
\setcounter{chapter}{-2}          % so that \chapter increments to -1
\stepcounter{chapter}
\chapter*{Lecture $-1$: Motivation}
\addcontentsline{toc}{chapter}{Lecture $-1$: Motivation}
\markboth{Lecture $-1$: Motivation}{}

\section*{Why Study Tensors?}

The principal motivation for studying tensors is \textbf{geometry}---not the elementary
geometry of basic shapes and trigonometry, but the rich, often counter-intuitive
geometry that arises in modern physics and mathematics.

\subsection*{General Relativity}

In Einstein's general theory of relativity, space-time is \emph{curved}, and light
can travel along curved paths near massive bodies such as stars and black holes.
The universe itself is expanding: beginning from the Big Bang, space has been
growing ever since.

To make these ideas mathematically precise we need tensors.  Einstein's field
equations,
\begin{equation}
    R_{\mu\nu} - \tfrac{1}{2}\,R\,g_{\mu\nu} + \Lambda\,g_{\mu\nu}
    = \frac{8\pi G}{c^{4}}\,T_{\mu\nu},
\end{equation}
actually encode \textbf{16 coupled equations}, and every symbol that appears
($R_{\mu\nu}$, $g_{\mu\nu}$, $T_{\mu\nu}$, etc.)\ is a tensor.  The most
important tensor in general relativity is the \textbf{metric tensor}~$g_{\mu\nu}$,
a $4\times 4$ rank-two tensor that lets us measure lengths and angles in the curved
geometry of space-time.

\subsection*{Quantum Mechanics and Quantum Computing}

In quantum mechanics, physical states are \emph{vectors}.
\begin{itemize}[nosep]
    \item \textbf{Quantum superposition} is simply a \emph{linear combination} of
          state vectors.
    \item \textbf{Quantum entanglement} arises when two quantum systems are combined
          via the \textbf{tensor product}~$\otimes$.  The tensor product takes the
          state spaces of two individual systems and merges them into a single,
          larger space in which the entangled system lives.
\end{itemize}

\subsection*{Prerequisites}

These notes assume basic familiarity with linear algebra: matrix multiplication,
linear combinations, and the dot product.  Calculus is \emph{not} required; tensors
can be developed entirely within the framework of linear algebra.

% ======================================================================
%  CHAPTER 0 : TENSOR DEFINITION
% ======================================================================
\chapter{Tensor Definition}\label{ch:tensor-def}

\section{The Array Definition (First Attempt)}

\begin{quote}
    \emph{A tensor is a multi-dimensional array of numbers.}
\end{quote}

Under this view, tensors are grids of numbers of increasing dimension:
\begin{center}
\renewcommand{\arraystretch}{1.3}
\begin{tabular}{lcc}
    \toprule
    \textbf{Name} & \textbf{Rank} & \textbf{Example} \\
    \midrule
    Scalar   & 0 & $\pi,\; 42,\; \tfrac{3}{4}$ \\
    Vector   & 1 & $\begin{pmatrix} v_1 \\ v_2 \\ v_3 \end{pmatrix}$ \\
    Matrix   & 2 & $\begin{pmatrix} a & b \\ c & d \end{pmatrix}$ \\
    Rank-3 tensor & 3 & (a three-dimensional ``cube'' of numbers) \\
    \bottomrule
\end{tabular}
\end{center}

\begin{warning}
This definition is \textbf{incomplete}.  Tensors can be \emph{represented} as
multi-dimensional arrays, but they carry genuine \textbf{geometric meaning} that a
bare array of numbers does not capture.
\end{warning}

\section{The Coordinate Definition (Second Attempt)}

\begin{quote}
    \emph{A tensor is an object that is \textbf{invariant} under a change of
    coordinates, but whose \textbf{components} change in a special, predictable way
    under a change of coordinates.}
\end{quote}

\paragraph{Invariance.}
Consider a pencil pointing toward a door.  Its length and orientation are
\emph{intrinsic} properties that do not depend on any coordinate system.

\paragraph{Components depend on coordinates.}
If we introduce a coordinate system, we can decompose the pencil into components
along each basis direction.  Changing the coordinate system changes the components
but \emph{not} the pencil itself.

The ``special, predictable way'' in which components change is governed by
\textbf{forward} and \textbf{backward} transformations, which we develop in the next
chapter.

\section{The Abstract Definition (Best Definition)}\label{sec:abstract-def}

\begin{definition}[Tensor --- abstract definition]
A \textbf{tensor} is a collection of \textbf{vectors} and \textbf{covectors}
combined together using the \textbf{tensor product}.
\end{definition}

This definition is short and elegant but requires us to understand three ingredients:
\emph{vectors}, \emph{covectors}, and the \emph{tensor product}.  The remainder of
these notes builds up each ingredient in turn.

% ======================================================================
%  CHAPTER 1 : FORWARD AND BACKWARD TRANSFORMATIONS
% ======================================================================
\chapter{Forward and Backward Transformations}\label{ch:fwd-bwd}

\section{Setup: Old and New Bases}

Let $V$ be an $n$-dimensional vector space with two bases:
\begin{itemize}[nosep]
    \item \textbf{Old basis:} $\{e_1, e_2, \dots, e_n\}$.
    \item \textbf{New basis:} $\{\tildee{1}, \tildee{2}, \dots, \tildee{n}\}$.
\end{itemize}

\section{The Forward Transformation}

The \textbf{forward transformation} builds the new basis vectors from the old:
\begin{equation}\label{eq:fwd}
    \boxed{\tildee{i} = \sum_{j=1}^{n} F^{j}{}_{i}\, e_j}
\end{equation}
The coefficients $F^{j}{}_{i}$ are stored in an $n\times n$ matrix~$F$, called the
\textbf{forward matrix}.

\begin{remark}
Throughout these notes, basis vectors are written as a \emph{row} that multiplies the
transformation matrix from the left:
$\bigl(\tildee{1}\;\; \tildee{2}\bigr)
= \bigl(e_1\;\; e_2\bigr)\,F.$
Consequently the entries of $F$ are \emph{transposed} relative to the scalar
coefficients in the corresponding linear-combination equations.
\end{remark}

\section{The Backward Transformation}

The \textbf{backward transformation} reverses the process, building old basis vectors
from the new:
\begin{equation}\label{eq:bwd}
    \boxed{e_i = \sum_{j=1}^{n} B^{j}{}_{i}\, \tildee{j}}
\end{equation}
where $B$ is the \textbf{backward matrix}.

\section{$F$ and $B$ Are Inverses}

\begin{proposition}
$F$ and $B$ are inverse matrices:
\[
    FB = BF = I, \qquad\text{i.e.}\quad B = F^{-1}.
\]
Equivalently, in component form,
\begin{equation}\label{eq:fb-inv}
    \sum_{j} F^{j}{}_{i}\, B^{k}{}_{j} = \delta^{k}_{i},
\end{equation}
where $\delta^{k}_{i}$ is the \textbf{Kronecker delta}:
\[
    \delta^{k}_{i} = \begin{cases} 1 & \text{if } i = k,\\ 0 & \text{if } i \neq k. \end{cases}
\]
\end{proposition}

\begin{proof}
Substitute the forward transformation~\eqref{eq:fwd} into the backward
transformation~\eqref{eq:bwd} (or vice-versa).  Requiring that $e_i = e_i$ forces
the product of the two matrices to equal the identity.
\end{proof}

\begin{remark}
We use separate letters $F$ and $B$ (rather than $F$ and $F^{-1}$) purely for visual
clarity.
\end{remark}

% ======================================================================
%  CHAPTER 2 : VECTORS
% ======================================================================
\chapter{Vectors}\label{ch:vectors}

\section{What Is a Vector?}

Three progressively better definitions:

\paragraph{Definition 1 --- Array.}
A vector is a one-dimensional array (list) of numbers.
This is convenient for computation but conflates a vector with its
\emph{components}, which depend on a coordinate system.

\paragraph{Definition 2 --- Arrow.}
A vector is an arrow in space.  Arrows can be added (tip-to-tail) and scaled
(stretched or reversed).  This is geometrically intuitive but limited to
\emph{Euclidean} vectors; other vector types exist.

\paragraph{Definition 3 --- Vector space (best).}
\begin{definition}[Vector space]
A \textbf{vector space} is a collection $(V, S, +, \cdot\,)$ where
\begin{enumerate}[nosep]
    \item $V$ is a set of \textbf{vectors},
    \item $S$ is a set of \textbf{scalars} (typically $\R$),
    \item $+$ is an \textbf{addition} rule combining two vectors,
    \item $\cdot$ is a \textbf{scalar multiplication} rule.
\end{enumerate}
A \textbf{vector} is any element of~$V$.
\end{definition}

\section{Vector Components and Change of Basis}

Given a basis $\{e_i\}$, any vector $v\in V$ can be written
\[
    v = \sum_{i} v^{i}\, e_i,
\]
where $v^{i}$ are the \textbf{components} of $v$ in that basis.

\begin{example}[Two-dimensional]
In the old basis $\{e_1, e_2\}$ a vector may have components
$v^{i}_{(e)} = \bigl(1,\;1.5\bigr)$, while in the new basis
$\{\tildee{1},\tildee{2}\}$ the same vector has components
$\tilde{v}^{i} = \bigl(1,\;2\bigr)$.
\end{example}

\subsection*{Transformation of vector components}

Crucially, vector components transform \textbf{opposite} to basis vectors:

\begin{equation}\label{eq:vec-comp-fwd}
    \boxed{\tilde{v}^{i} = \sum_{j} B^{i}{}_{j}\, v^{j}}
    \qquad\text{(old $\to$ new: use $B$)}
\end{equation}
\begin{equation}\label{eq:vec-comp-bwd}
    \boxed{v^{i} = \sum_{j} F^{i}{}_{j}\, \tilde{v}^{j}}
    \qquad\text{(new $\to$ old: use $F$)}
\end{equation}

\section{Contravariance of Vector Components}\label{sec:contravariance}

\begin{itemize}[nosep]
    \item \textbf{Basis vectors} go from old to new via $F$ (forward).
    \item \textbf{Vector components} go from old to new via $B$ (backward).
\end{itemize}
They transform in \emph{opposite} ways.  This is intuitively clear: if we double
the basis vectors, the components halve (the measuring sticks are now bigger, so
fewer are needed).  If we rotate the basis one way, the components effectively
rotate the other way.

Because vector components behave \textbf{contrary} to the basis, they are called
\textbf{contravariant}.

\begin{notation}
We write vector component indices \emph{upstairs}: $v^{i}$.  This is a reminder
that they are contravariant (opposite to basis indices, which sit downstairs).
Upper indices should \emph{not} be confused with exponents; exponents appear very
rarely in tensor analysis.
\end{notation}

% ======================================================================
%  CHAPTER 3 : COVECTORS
% ======================================================================
\chapter{Covectors}\label{ch:covectors}

\section{What Is a Covector?}

\begin{definition}[Covector]
A \textbf{covector} (also called a \emph{linear form} or \emph{one-form}) is a
\textbf{linear function}
\[
    \alpha : V \longrightarrow \R
\]
that maps a vector to a scalar and satisfies the \textbf{linearity} properties:
\begin{align}
    \alpha(n\,v) &= n\,\alpha(v), \label{eq:cov-scale}\\
    \alpha(v + w) &= \alpha(v) + \alpha(w). \label{eq:cov-add}
\end{align}
\end{definition}

As a first approximation, covectors correspond to \textbf{row vectors} in a given
basis.  However, row vectors and column vectors are fundamentally different objects;
they only coincide in an orthonormal basis.

\section{Visualizing Covectors as Stacks of Lines}

A covector in 2D can be thought of as a function $\alpha_1 x + \alpha_2 y$.  Its
level sets (lines of constant value) form an oriented family of equally spaced,
parallel lines.

\begin{itemize}[nosep]
    \item The \textbf{value} $\alpha(v)$ equals the number of lines that the arrow
          $v$ \emph{pierces}.
    \item Piercing in the direction of increasing value gives a \emph{positive}
          result; against the orientation gives a \emph{negative} result.
    \item An arrow parallel to the lines pierces none, giving $\alpha(v) = 0$.
\end{itemize}

\subsection*{Scaling and adding covectors}

\begin{itemize}[nosep]
    \item \textbf{Scaling} a covector by $n$ makes the stack $n$ times \emph{denser}
          (not wider).
    \item \textbf{Adding} two covectors $\beta + \gamma$ produces a new stack whose
          density in each direction combines the densities of $\beta$ and $\gamma$.
\end{itemize}

\section{The Dual Space $V^*$}

Since covectors can be added and scaled in a meaningful way, they form a vector
space of their own.

\begin{definition}[Dual space]
The set of all covectors acting on a vector space $V$ is called the
\textbf{dual space}, denoted $V^{*}$.  It has its own addition and scalar
multiplication rules (inherited from the linearity of the covectors).
\end{definition}

\section{Covector Components and the Dual Basis}\label{sec:dual-basis}

Given a basis $\{e_1, \dots, e_n\}$ for $V$, we define the \textbf{dual basis}
$\{\epsilon^1, \dots, \epsilon^n\}$ for $V^*$ by
\begin{equation}\label{eq:dual-basis-def}
    \boxed{\epsilon^{i}(e_j) = \delta^{i}_{j}.}
\end{equation}
The dual basis covectors \textbf{project out} vector components:
\[
    \epsilon^{i}(v) = v^{i}.
\]
Any covector $\alpha$ can be expanded as
\[
    \alpha = \sum_{i} \alpha_i\, \epsilon^{i},
\]
where the components are $\alpha_i = \alpha(e_i)$---the value of $\alpha$ on the
$i$-th basis vector.

\section{Covector Component Transformation Rules}

Covector components transform using the \textbf{forward} matrix when going from old
to new (the \emph{same} direction as basis vectors):
\begin{equation}\label{eq:cov-comp-fwd}
    \boxed{\tilde{\alpha}_i = \sum_{j} F^{j}{}_{i}\, \alpha_j}
    \qquad\text{(old $\to$ new: use $F$)}
\end{equation}
\begin{equation}\label{eq:cov-comp-bwd}
    \boxed{\alpha_i = \sum_{j} B^{j}{}_{i}\, \tilde{\alpha}_j}
    \qquad\text{(new $\to$ old: use $B$)}
\end{equation}

\begin{remark}[Intuition]
When the basis vectors grow by a factor of 2, each basis vector now pierces
\emph{more} covector lines, so the covector components \emph{also} grow.
Covector components and basis vectors transform in the \emph{same} direction.
\end{remark}

Because covector components transform in the same way as basis vectors, they are
called \textbf{covariant}, and their indices are written \emph{downstairs}:
$\alpha_{i}$.

\section{Summary of Transformation Rules So Far}

\begin{center}
\renewcommand{\arraystretch}{1.4}
\begin{tabular}{lcc}
    \toprule
    \textbf{Object} & \textbf{Old $\to$ New} & \textbf{Type} \\
    \midrule
    Basis vectors $e_i$            & $F$ (forward)  & covariant \\
    Vector components $v^{i}$      & $B$ (backward) & contravariant \\
    Basis covectors $\epsilon^{i}$ & $B$ (backward) & contravariant \\
    Covector components $\alpha_i$ & $F$ (forward)  & covariant \\
    \bottomrule
\end{tabular}
\end{center}

% ======================================================================
%  CHAPTER 4 : LINEAR MAPS
% ======================================================================
\chapter{Linear Maps}\label{ch:linear-maps}

\section{Three Definitions of a Linear Map}

\paragraph{Coordinate definition.}
A linear map is represented by a \textbf{matrix}---a 2D array of numbers.  The
$i$-th column of the matrix tells us where a copy of the $i$-th basis vector is
sent.

\paragraph{Geometric definition.}
A linear map is a spatial transformation that
\begin{enumerate}[nosep]
    \item keeps grid lines \textbf{parallel},
    \item keeps grid lines \textbf{evenly spaced},
    \item keeps the \textbf{origin fixed}.
\end{enumerate}
Examples include stretches, rotations, and shears.  Translations are
\emph{not} linear maps because they move the origin.

\paragraph{Abstract definition (best).}
\begin{definition}[Linear map]
A \textbf{linear map} is a function $L : V \to W$ satisfying
\begin{align}
    L(v + w) &= L(v) + L(w), \label{eq:lm-add}\\
    L(a\,v)  &= a\,L(v). \label{eq:lm-scale}
\end{align}
\end{definition}
Covectors are a special case: they are linear maps $V \to \R$.  The only
difference is that covectors output a \emph{scalar}, while general linear maps
output a \emph{vector}.

\section{Deriving Matrix Multiplication from Linearity}

Let $L : V \to V$ be a linear map.  Expanding an input vector $v$ in the basis:
\[
    L(v) = L\!\Bigl(\sum_j v^{j}\,e_j\Bigr)
         = \sum_j v^{j}\,L(e_j).
\]
Define the coefficients $L^{i}{}_{j}$ by
\[
    L(e_j) = \sum_i L^{i}{}_{j}\, e_i.
\]
Then the output vector $w = L(v)$ has components
\begin{equation}\label{eq:mat-mult}
    \boxed{w^{i} = \sum_j L^{i}{}_{j}\, v^{j},}
\end{equation}
which is exactly the standard formula for matrix--vector multiplication.

\section{Transformation Rules for Linear Map Components}\label{sec:lm-transform}

To move the matrix components from the old basis to the new:
\begin{equation}\label{eq:lm-fwd}
    \boxed{\tilde{L}^{i}{}_{j}
    = \sum_{k,\ell} B^{i}{}_{k}\, L^{k}{}_{\ell}\, F^{\ell}{}_{j}}
    \qquad\Longleftrightarrow\qquad
    \tilde{L} = B\,L\,F.
\end{equation}
The reverse is
\begin{equation}\label{eq:lm-bwd}
    L^{i}{}_{j}
    = \sum_{k,\ell} F^{i}{}_{k}\, \tilde{L}^{k}{}_{\ell}\, B^{\ell}{}_{j}
    \qquad\Longleftrightarrow\qquad
    L = F\,\tilde{L}\,B.
\end{equation}

\begin{remark}[Interpretation]
The path through the transformation diagram is:
\[
    \text{new components}
    \xrightarrow{F} \text{old components}
    \xrightarrow{L} \text{old output components}
    \xrightarrow{B} \text{new output components}.
\]
\end{remark}

\section{Einstein Summation Convention}\label{sec:einstein}

Whenever an index appears \textbf{once upstairs and once downstairs} in the same
term, summation over that index is implied.  We drop the $\sum$ sign entirely.

\begin{example}
The matrix transformation rule becomes simply
\[
    \tilde{L}^{i}{}_{j} = B^{i}{}_{k}\, L^{k}{}_{\ell}\, F^{\ell}{}_{j}.
\]
\end{example}

\begin{remark}[Kronecker delta as index canceller]
Since $M^{i}{}_{k}\,\delta^{k}_{j} = M^{i}{}_{j}$, the Kronecker delta effectively
``cancels'' the summed index and replaces it with the free one.
\end{remark}

\section{Tensor Type of a Linear Map}

Linear maps transform with \emph{one} backward (contravariant) and \emph{one}
forward (covariant) rule.  They are therefore \textbf{$(1,1)$-tensors}.

% ======================================================================
%  CHAPTER 5 : THE METRIC TENSOR
% ======================================================================
\chapter{The Metric Tensor}\label{ch:metric}

\section{Measuring Lengths: Beyond Pythagoras}

In an orthonormal basis, $|v|^{2} = (v^{1})^{2} + (v^{2})^{2} + \cdots$
(Pythagoras).  In a \emph{general} (non-orthonormal) basis this formula
\textbf{fails}: the coordinate triangle need not be right-angled.

The correct formula uses the \textbf{dot product}:
\begin{equation}\label{eq:len-dot}
    |v|^{2} = v \cdot v.
\end{equation}
Expanding in a basis:
\begin{equation}\label{eq:len-expand}
    |v|^{2}
    = \Bigl(\sum_i v^{i}\,e_i\Bigr)\cdot\Bigl(\sum_j v^{j}\,e_j\Bigr)
    = \sum_{i,j} v^{i}\,v^{j}\,(e_i \cdot e_j).
\end{equation}
The dot products of basis vectors are the key quantities.

\section{Definition of the Metric Tensor}

\begin{definition}[Metric tensor]
The \textbf{metric tensor} $g$ has components
\begin{equation}\label{eq:metric-def}
    \boxed{g_{ij} = e_i \cdot e_j.}
\end{equation}
\end{definition}

Since the dot product is symmetric ($e_i \cdot e_j = e_j \cdot e_i$), the metric
tensor is a \textbf{symmetric matrix}: $g_{ij} = g_{ji}$.

The squared length of a vector is then:
\begin{equation}\label{eq:metric-length}
    \boxed{|v|^{2} = \sum_{i,j} g_{ij}\, v^{i}\, v^{j} = g_{ij}\, v^{i}\, v^{j}.}
\end{equation}

\begin{example}
In an orthonormal basis, $g_{ij} = \delta_{ij}$ (the identity matrix), and the
formula reduces to Pythagoras: $|v|^{2} = (v^{1})^{2} + \cdots + (v^{n})^{2}$.
\end{example}

\section{Measuring Angles}

The dot product of two vectors is
\[
    v \cdot w = g_{ij}\, v^{i}\, w^{j},
\]
and the angle $\theta$ between $v$ and $w$ satisfies
\begin{equation}\label{eq:angle}
    \boxed{\cos\theta = \frac{v \cdot w}{|v|\;|w|}
    = \frac{g_{ij}\, v^{i}\, w^{j}}
           {\sqrt{g_{ij}\, v^{i}\, v^{j}}\;\sqrt{g_{kl}\, w^{k}\, w^{l}}}.}
\end{equation}

\section{Transformation Rules for the Metric Tensor}

The metric tensor components transform with \textbf{two forward} matrices:
\begin{equation}\label{eq:metric-fwd}
    \boxed{\tilde{g}_{ij} = F^{k}{}_{i}\, F^{\ell}{}_{j}\, g_{k\ell}}
    \qquad\text{(old $\to$ new)}
\end{equation}
\begin{equation}\label{eq:metric-bwd}
    g_{ij} = B^{k}{}_{i}\, B^{\ell}{}_{j}\, \tilde{g}_{k\ell}
    \qquad\text{(new $\to$ old)}
\end{equation}
Because it uses \emph{two covariant} rules, the metric tensor is a
\textbf{$(0,2)$-tensor}.

\section{Invariance of Length}

Using the transformation rules for $g_{ij}$ and $v^{i}$:
\[
    \tilde{g}_{ij}\,\tilde{v}^{i}\,\tilde{v}^{j}
    = \bigl(F^{k}{}_{i}\,F^{\ell}{}_{j}\,g_{k\ell}\bigr)
      \bigl(B^{i}{}_{m}\,v^{m}\bigr)
      \bigl(B^{j}{}_{n}\,v^{n}\bigr)
    = g_{mn}\,v^{m}\,v^{n},
\]
since $F^{k}{}_{i}\,B^{i}{}_{m} = \delta^{k}_{m}$.  The squared length is the
same in every coordinate system, confirming that length is \textbf{invariant}.

% ======================================================================
%  CHAPTER 6 : BILINEAR FORMS
% ======================================================================
\chapter{Bilinear Forms}\label{ch:bilinear}

\section{Definition}

\begin{definition}[Bilinear form]
A \textbf{bilinear form} is a function $B : V \times V \to \R$ that is
\textbf{linear in each input separately}:
\begin{align}
    B(a\,v, w) &= a\,B(v, w) = B(v, a\,w), \label{eq:bf-scale}\\
    B(v_1 + v_2,\, w) &= B(v_1, w) + B(v_2, w), \label{eq:bf-add1}\\
    B(v,\, w_1 + w_2) &= B(v, w_1) + B(v, w_2). \label{eq:bf-add2}
\end{align}
\end{definition}

In a given basis, the output is computed by
\begin{equation}
    B(v,w) = B_{ij}\, v^{i}\, w^{j}.
\end{equation}

\begin{warning}
Do \textbf{not} scale \emph{both} inputs simultaneously:
$B(a\,v,\,a\,w) = a^{2}\,B(v,w) \neq a\,B(v,w)$
(unless $a = 0$ or $1$).  Likewise, do not add both inputs at once without
distributing carefully---four terms will result.
\end{warning}

\section{Bilinear Forms vs.\ the Metric Tensor}

The metric tensor is a \emph{special} bilinear form with two extra properties:
\begin{enumerate}[nosep]
    \item \textbf{Symmetry:} $g(v,w) = g(w,v)$, i.e.\ $g_{ij} = g_{ji}$.
    \item \textbf{Positive-definiteness:} $g(v,v) \geq 0$ for all $v$, with
          equality only when $v = 0$.
\end{enumerate}
Not every bilinear form satisfies these conditions.  Bilinear forms that do are
valid metric tensors; those that do not are still $(0,2)$-tensors, but they are
\emph{not} metrics.

\section{Terminology}

\begin{itemize}[nosep]
    \item A \textbf{form} is a function that takes vectors as input and outputs a
          number.
    \item A covector is a \textbf{linear form} (or \textbf{one-form}): one input,
          linear.
    \item A bilinear form is a \textbf{two-form}: two inputs, linear in each
          separately.
\end{itemize}

% ======================================================================
%  CHAPTER 7 : THE GENERAL TRANSFORMATION LAW
% ======================================================================
\chapter{The General Tensor Transformation Law}\label{ch:general-law}

Collecting the results from previous chapters, any object whose components carry $m$
upper (contravariant) indices and $n$ lower (covariant) indices transforms as:

\begin{equation}\label{eq:general-transform}
    \boxed{
    \tilde{T}^{i_1 \cdots i_m}{}_{j_1 \cdots j_n}
    = B^{i_1}{}_{k_1} \cdots B^{i_m}{}_{k_m}\;
      F^{\ell_1}{}_{j_1} \cdots F^{\ell_n}{}_{j_n}\;
      T^{k_1 \cdots k_m}{}_{\ell_1 \cdots \ell_n}
    }
\end{equation}

\begin{itemize}[nosep]
    \item Each \textbf{upper index} transforms with a $B$ (backward, contravariant).
    \item Each \textbf{lower index} transforms with an $F$ (forward, covariant).
\end{itemize}

\begin{definition}[Tensor type]
A tensor with $m$ contravariant (upper) indices and $n$ covariant (lower) indices is
called an $\bm{(m,n)}$\textbf{-tensor}.
\end{definition}

\begin{center}
\renewcommand{\arraystretch}{1.4}
\begin{tabular}{lcl}
    \toprule
    \textbf{Tensor} & \textbf{Type} & \textbf{Transforms with} \\
    \midrule
    Scalar           & $(0,0)$ & (invariant) \\
    Vector component $v^{i}$ & $(1,0)$ & one $B$ \\
    Covector component $\alpha_i$ & $(0,1)$ & one $F$ \\
    Linear map $L^{i}{}_{j}$ & $(1,1)$ & one $B$, one $F$ \\
    Metric tensor $g_{ij}$   & $(0,2)$ & two $F$'s \\
    \bottomrule
\end{tabular}
\end{center}

\textbf{Anything that follows the rule~\eqref{eq:general-transform} under a change
of coordinates is a tensor.}

% ======================================================================
%  CHAPTER 8 : THE TENSOR PRODUCT — BUILDING TENSORS
% ======================================================================
\chapter{The Tensor Product}\label{ch:tensor-product}

\section{Vector--Covector Pairs as Linear Maps}\label{sec:vec-cov-pairs}

Multiplying a column vector (left) by a row vector (right) produces a matrix:
\[
    \begin{pmatrix} a \\ b \end{pmatrix}
    \begin{pmatrix} c & d \end{pmatrix}
    = \begin{pmatrix} ac & ad \\ bc & bd \end{pmatrix}.
\]
In tensor notation, writing a vector $e_i$ next to a covector $\epsilon^{j}$ gives
a basis linear map $e_i \otimes \epsilon^{j}$.  Any linear map $L$ can be written
\begin{equation}\label{eq:L-tensor-expand}
    L = L^{i}{}_{j}\; e_i \otimes \epsilon^{j}.
\end{equation}

\subsection*{Pure vs.\ impure matrices}

\begin{itemize}[nosep]
    \item A \textbf{pure} matrix can be factored into a single column--row product.
          All its columns are scalar multiples of each other, so the linear map
          sends every input to the same direction.
    \item An \textbf{impure} matrix cannot be so factored.  It is a
          \emph{linear combination} of pure matrices and can produce outputs
          pointing in different directions.
\end{itemize}

The four products $e_i \otimes \epsilon^{j}$ ($i,j = 1,2$) form a
\textbf{basis for all $2\times 2$ linear maps}.

\subsection*{Action on a vector}

Let $L = L^{i}{}_{j}\; e_i \otimes \epsilon^{j}$ act on $v = v^{k}\,e_k$:
\[
    L(v) = L^{i}{}_{j}\,v^{k}\; e_i\;\underbrace{\epsilon^{j}(e_k)}_{\delta^{j}_{k}}
         = L^{i}{}_{j}\,v^{j}\; e_i,
\]
recovering the standard matrix--vector multiplication formula.

\section{Covector--Covector Pairs as Bilinear Forms}\label{sec:cov-cov-pairs}

Similarly, a bilinear form can be expanded as
\begin{equation}
    B = B_{ij}\; \epsilon^{i} \otimes \epsilon^{j}.
\end{equation}
Acting on two vectors $v = v^{k}\,e_k$ and $w = w^{\ell}\,e_\ell$:
\[
    B(v,w) = B_{ij}\,v^{k}\,w^{\ell}\;
    \underbrace{\epsilon^{i}(e_k)}_{\delta^{i}_k}
    \underbrace{\epsilon^{j}(e_\ell)}_{\delta^{j}_\ell}
    = B_{ij}\,v^{i}\,w^{j}.
\]

The array shape is a \textbf{row of rows}---both vector inputs can be written as
columns, and the multiplication works out naturally.

\section{Tensor Product vs.\ Kronecker Product}\label{sec:tp-vs-kp}

The \textbf{tensor product} $\otimes$ and the \textbf{Kronecker product} use the
same symbol but operate in different domains:

\begin{center}
\renewcommand{\arraystretch}{1.3}
\begin{tabular}{lll}
    \toprule
    & \textbf{Tensor product} & \textbf{Kronecker product} \\
    \midrule
    Operands & abstract tensors & arrays of numbers \\
    Result   & a new tensor     & a new array \\
    \bottomrule
\end{tabular}
\end{center}

\noindent They are closely related: the \emph{components} obtained from the tensor
product are exactly the \emph{entries} of the array obtained from the Kronecker
product.  The tensor product works in the land of algebraic symbols; the Kronecker
product works in the land of arrays.

\section{Benefits of the Tensor Product Viewpoint}

Given the expansion $L = L^{i}{}_{j}\;e_i \otimes \epsilon^{j}$:

\begin{enumerate}[nosep]
    \item \textbf{Transformation rules for free.}  Transform each basis vector and
          basis covector individually using the rules already known.
    \item \textbf{Multiplication formulas for free.}  Pass vector inputs to
          covectors and apply linearity + Kronecker delta.
    \item \textbf{Correct array shapes for free.}  The Kronecker product of the
          component arrays gives the correct multi-dimensional array.
\end{enumerate}

% ======================================================================
%  CHAPTER 9 : GENERAL TENSORS AND TENSOR PRODUCT SPACES
% ======================================================================
\chapter{General Tensors and Tensor Product Spaces}\label{ch:general-tensors}

\section{Arbitrary Vector--Covector Combinations}

The tensor product viewpoint extends to \emph{any} combination of vectors and
covectors.

\begin{example}\leavevmode
\begin{itemize}[nosep]
    \item A $(2,0)$-tensor: $D = D^{ab}\; e_a \otimes e_b$ \quad (vector--vector
          pair).
    \item A $(1,2)$-tensor: $Q = Q^{i}{}_{jk}\; e_i \otimes \epsilon^{j}
          \otimes \epsilon^{k}$ \quad (vector--covector--covector triple).
\end{itemize}
\end{example}

Transformation rules follow immediately by transforming each basis element
individually.

\subsection*{Multiple contraction possibilities}

For higher-type tensors, there is generally \emph{no unique} way to ``multiply''
two tensors---the Einstein component notation must specify \emph{which} indices are
summed.  For instance, $Q^{i}{}_{jk}\,D^{jk}$, $Q^{i}{}_{jk}\,D^{kj}$,
and $Q^{i}{}_{jk}\,D^{ja}$ are all different operations.

\section{Tensor Product Spaces}

The tensor product of vector spaces $V$ and $W$ produces a new vector space
$V \otimes W$ whose elements obey the \textbf{scaling rule}
\[
    n\,(a \otimes b) = (n\,a) \otimes b = a \otimes (n\,b),
\]
and the \textbf{addition rules}
\begin{align}
    (a_1 + a_2) \otimes b &= a_1 \otimes b + a_2 \otimes b, \\
    a \otimes (b_1 + b_2) &= a \otimes b_1 + a \otimes b_2.
\end{align}

\begin{example}[Tensor product spaces]\leavevmode
\begin{itemize}[nosep]
    \item $(1,1)$-tensors live in $V \otimes V^*$.
    \item $(0,2)$-tensors live in $V^* \otimes V^*$.
    \item $(2,0)$-tensors live in $V \otimes V$.
    \item $(1,2)$-tensors live in $V \otimes V^* \otimes V^*$.
\end{itemize}
\end{example}

We can build arbitrarily large tensor product spaces:
\[
    \underbrace{V \otimes \cdots \otimes V}_{m} \otimes
    \underbrace{V^* \otimes \cdots \otimes V^*}_{n}
\]
contains all $(m,n)$-tensors.  Component indices from $V$ go \emph{upstairs};
indices from $V^*$ go \emph{downstairs}.

\section{Tensors as Multilinear Maps}

\begin{definition}[Multilinear map]
A \textbf{multilinear map} is a function of several variables that is
\textbf{linear in each variable separately} (while all others are held fixed):
\begin{align}
    T(\dots, n\,v_k, \dots) &= n\,T(\dots, v_k, \dots), \\
    T(\dots, v_k + w_k, \dots) &= T(\dots, v_k, \dots) + T(\dots, w_k, \dots).
\end{align}
\end{definition}

All tensors, when used as functions, are multilinear maps.  This is a direct
consequence of the linearity of each covector and vector factor in the tensor
product.

% ======================================================================
%  CHAPTER 10 : RAISING AND LOWERING INDICES
% ======================================================================
\chapter{Raising and Lowering Indices}\label{ch:raise-lower}

\section{Creating Vector--Covector Partners}

We seek a natural correspondence between vectors in $V$ and covectors in $V^*$.

\paragraph{Na\"ive attempt.}
Pair basis vector $e_i$ with basis covector $\epsilon^{i}$.  This looks clean in
one basis but breaks down under a change of basis: basis vectors transform
covariantly while basis covectors transform contravariantly, so the pairing becomes
inconsistent.

\paragraph{Correct approach.}
Given a vector $v \in V$, define the covector
\begin{equation}\label{eq:v-flat}
    v^{\flat} \;:\; w \longmapsto v \cdot w.
\end{equation}
This is indeed a covector (an element of $V^*$) because:
\begin{enumerate}[nosep]
    \item It maps vectors to scalars.
    \item It is linear: $v \cdot (n\,w) = n\,(v \cdot w)$ and
          $v \cdot (w_1 + w_2) = v \cdot w_1 + v \cdot w_2$.
\end{enumerate}

This correspondence does \textbf{not} depend on any basis and behaves well under
changes of coordinates: scaling $v$ by 2 also scales $v^{\flat}$ by 2.

\section{Lowering Indices with the Metric Tensor}

The components of $v^{\flat}$ are obtained via the metric tensor:
\begin{equation}\label{eq:lower}
    \boxed{v_i = g_{ij}\, v^{j}.}
\end{equation}
Here $v^{j}$ are the (contravariant) vector components and $v_i$ are the
(covariant) covector components of the partner covector $v^{\flat}$.

\begin{warning}
$v_i$ and $v^{i}$ are \textbf{not the same thing}.  Converting between them
\textbf{requires} the metric tensor.  They are only equal in the special case of an
orthonormal basis where $g_{ij} = \delta_{ij}$.
\end{warning}

\section{The Inverse Metric Tensor}

\begin{definition}[Inverse metric tensor]
The \textbf{inverse metric tensor} $g^{ij}$ is defined by
\begin{equation}\label{eq:inv-metric}
    \boxed{g^{ik}\, g_{kj} = \delta^{i}_{j}.}
\end{equation}
It is a $(2,0)$-tensor living in $V \otimes V$.
\end{definition}

\section{Raising Indices}

Multiplying both sides of the lowering equation~\eqref{eq:lower} by $g^{ik}$:
\begin{equation}\label{eq:raise}
    \boxed{v^{i} = g^{ij}\, v_j.}
\end{equation}

\begin{center}
\renewcommand{\arraystretch}{1.3}
\begin{tabular}{lcc}
    \toprule
    \textbf{Operation} & \textbf{Formula} & \textbf{Uses} \\
    \midrule
    Lower an index & $v_i = g_{ij}\,v^{j}$ & metric tensor $g_{ij}$ \\
    Raise an index & $v^{i} = g^{ij}\,v_j$  & inverse metric $g^{ij}$ \\
    \bottomrule
\end{tabular}
\end{center}

\section{Raising and Lowering on General Tensors}

These operations extend to tensors of any type.  For example, given a $(1,2)$-tensor
$Q^{i}{}_{jk} \in V \otimes V^* \otimes V^*$, we can raise the $j$-index:
\[
    Q'^{i\ell}{}_{k} = g^{\ell j}\, Q^{i}{}_{jk},
\]
producing a $(2,1)$-tensor in $V \otimes V \otimes V^*$.

Similarly, for a $(2,0)$-tensor $D^{ab}$, we can lower the $a$-index:
\[
    D'_{a}{}^{b} = g_{ac}\, D^{cb},
\]
producing a $(1,1)$-tensor.

\section{The Flat and Sharp Operators}

Musical notation provides a concise shorthand:

\begin{itemize}[nosep]
    \item $v^{\flat}$ (\textbf{flat}): converts a vector $v$ to its partner
          covector.  The flat symbol \emph{lowers} the index, just as a flat in
          music lowers a note's pitch.  Geometrically, it ``flattens'' a pointy
          arrow into a flat stack of lines.
    \item $\alpha^{\sharp}$ (\textbf{sharp}): converts a covector $\alpha$ to its
          partner vector.  The sharp symbol \emph{raises} the index, just as a
          sharp in music raises a note's pitch.  Geometrically, it ``sharpens'' a
          flat stack into a pointy arrow.
\end{itemize}

In formulas:
\begin{align}
    (v^{\flat})_i &= g_{ij}\, v^{j}, \\
    (\alpha^{\sharp})^{i} &= g^{ij}\, \alpha_j.
\end{align}

% ======================================================================
%  APPENDIX: SUMMARY
% ======================================================================
\appendix
\chapter{Quick Reference}\label{app:summary}

\section{Basis Transformations}

\begin{align}
    \tildee{i} &= F^{j}{}_{i}\, e_j
        &\quad& \text{(forward: old basis $\to$ new basis)} \\
    e_i &= B^{j}{}_{i}\, \tildee{j}
        &\quad& \text{(backward: new basis $\to$ old basis)}
\end{align}

\section{Component Transformations (Old $\to$ New)}

\begin{center}
\renewcommand{\arraystretch}{1.4}
\begin{tabular}{lccc}
    \toprule
    \textbf{Object} & \textbf{Type} & \textbf{Rule} & \textbf{Name} \\
    \midrule
    Vector components $v^{i}$    & $(1,0)$ & $\tilde{v}^{i} = B^{i}{}_{j}\,v^{j}$
        & contravariant \\
    Covector components $\alpha_i$ & $(0,1)$ & $\tilde{\alpha}_i = F^{j}{}_{i}\,\alpha_j$
        & covariant \\
    Linear map $L^{i}{}_{j}$     & $(1,1)$ & $\tilde{L}^{i}{}_{j} = B^{i}{}_{k}\,L^{k}{}_{\ell}\,F^{\ell}{}_{j}$
        & mixed \\
    Metric tensor $g_{ij}$       & $(0,2)$ & $\tilde{g}_{ij} = F^{k}{}_{i}\,F^{\ell}{}_{j}\,g_{k\ell}$
        & covariant \\
    \bottomrule
\end{tabular}
\end{center}

\section{Key Identities}

\begin{align}
    F^{j}{}_{i}\, B^{k}{}_{j} &= \delta^{k}_{i}
        &\quad& \text{($F$ and $B$ are inverses)} \\
    \epsilon^{i}(e_j) &= \delta^{i}_{j}
        &\quad& \text{(dual basis definition)} \\
    g_{ij}\, v^{j} &= v_i
        &\quad& \text{(index lowering)} \\
    g^{ij}\, v_j &= v^{i}
        &\quad& \text{(index raising)} \\
    g^{ik}\, g_{kj} &= \delta^{i}_{j}
        &\quad& \text{(inverse metric definition)}
\end{align}

\section{Tensor Product}

\begin{itemize}[nosep]
    \item \textbf{Scaling:}
          $n\,(a \otimes b) = (n\,a) \otimes b = a \otimes (n\,b)$.
    \item \textbf{Addition (left):}
          $(a_1 + a_2) \otimes b = a_1 \otimes b + a_2 \otimes b$.
    \item \textbf{Addition (right):}
          $a \otimes (b_1 + b_2) = a \otimes b_1 + a \otimes b_2$.
\end{itemize}

\section{Flat and Sharp Operators}

\begin{align}
    v^{\flat} &: w \mapsto v \cdot w
        &\quad& \text{(vector $\to$ covector, lowers index)} \\
    \alpha^{\sharp} &: \text{covector} \to \text{vector}
        &\quad& \text{(covector $\to$ vector, raises index)}
\end{align}

\end{document}
